\section{Generative AI: adelantar el control de riesgos a todo el ciclo de vida}

Lo que cambia la IA generativa no es solo que «se produzcan archivos nuevos»: también cambian la suplantación de identidad, las variantes de contenido, la integración en aplicaciones y la velocidad de distribución. Escanear únicamente en el punto final de carga a la plataforma traslada toda la responsabilidad al extremo de la cadena de suministro. Un enfoque más sólido es distribuir el threat model entre training data, model access, application, hosting/distribution e incident response, y verificarlo de forma continua en las tres etapas de development, deployment y maintenance.

\subsection{La threat surface abarca cuatro capas}

Esta sección parte de una perspectiva de cadena de suministro. Los desarrolladores del modelo controlan el entrenamiento y el fine-tuning; el servicio del modelo controla el access y el abuse monitoring; los desarrolladores de aplicaciones controlan el prompt flow, la user identity y los guardrail del producto; las plataformas de distribución controlan el hosting, el sharing, el reporting y la removal. Cada participante ve señales distintas, por lo que no se puede exigir que la plataforma del extremo recupere por sí sola la provenance o las relaciones de cuenta que ya se perdieron aguas arriba.

\lecturefigure{07-genai-threat-surface.png}{Los riesgos de la Generative AI abarcan las capas de desarrollo, acceso al modelo, aplicación y distribución.}{Subtítulos locales 05:45--10:20, 20:30--24:10; libro blanco Safety by Design de Thorn; redibujo conceptual.}

\begin{knowledgebox}{Lectura de la figura: escribir la capability y la obligation de cada capa}
La capa de development puede realizar dataset curation, evaluation y release gating; la capa de model access puede realizar authentication, rate limit, usage monitoring y abuse response; la capa de application puede diseñar una UX age-appropriate, contextual guardrail y user reporting; la capa de distribution puede ejecutar detection, removal, evidence preservation y lawful reporting. Ninguna capa posee todos los hechos, pero cada una debe conservar las señales mínimas de seguridad que puedan transmitirse a la capa siguiente.
\end{knowledgebox}

\begin{warningbox}{No convertir el threat model en un manual de ataque}
Los equipos de seguridad necesitan entender las categorías de abuso, los puntos de entrada y las señales observables, pero los apuntes públicos, los informes de pruebas y la documentación del producto no deben ofrecer pasos de evasión reproducibles, prompts sensibles ni flujos de generación. Una transparencia eficaz debe explicar la cobertura, los métodos de prueba, las categorías de falla, los tiempos de respuesta y los límites de gobernanza, en lugar de divulgar detalles que reduzcan de forma significativa el costo del abuso.
\end{warningbox}

\subsection{Safety by Design es una disciplina del ciclo de vida}

\term{Safety by Design} consiste en incorporar de forma continua el control de riesgos durante el diseño, el desarrollo, el despliegue y el mantenimiento, en lugar de añadir un moderation endpoint después del lanzamiento. Exige que el equipo escriba el threat model al definir los objetivos del producto, que realice evaluation/red teaming antes del entrenamiento y del release, que diseñe monitoring, escalation y user controls al desplegar, y que atienda el drift, los incident y el transparency reporting durante el mantenimiento.

\lecturefigure{08-safety-lifecycle.png}{Safety by Design abarca design, development, deployment y maintenance, no una única revisión de lanzamiento.}{Subtítulos locales 20:30--26:20; principios de Thorn/All Tech Is Human y NIST AI 600-1; redibujo conceptual.}

\begin{knowledgebox}{Lectura de la figura: el release gate es solo un punto intermedio}
La etapa de design define el intended use, el misuse y los grupos afectados; la etapa de development gobierna los datos, el modelo y las pruebas; la etapa de deployment configura el threshold, la human review, el appeal y el incident channel; la etapa de maintenance monitorea el drift, actualiza la policy, corrige vulnerabilidades y publica los avances. Un modelo que aprobó el launch checklist, si no cuenta con telemetry continua ni con responsables, puede fallar al cabo de unas semanas por cambios en el tráfico, el idioma o la integración con el producto.
\end{knowledgebox}

El \term{red teaming} (pruebas de equipo rojo) consiste en buscar activamente los modos de falla del sistema en un entorno autorizado y controlado; debe abarcar las salidas del modelo, los flujos de trabajo del producto, el abuso de cuentas, la cadena de suministro y el response process. Los resultados de las pruebas deben traducirse en owner, severity, fix, retest y residual risk, y no limitarse a mostrar un conjunto de capturas de pantalla de «inducción exitosa del modelo».

\teachervoice{Cordua describe los compromisos de la industria como un lenguaje común: los expertos en seguridad infantil no pueden reescribir el producto completo de cada empresa, pero sí pueden convertir los riesgos recurrentes en principles, standards y progress reporting que se puedan probar, de modo que la responsabilidad entre en el flujo normal de ingeniería.}

\subsection{Resumen de la sección}

La IA generativa requiere responsabilidad a lo largo de la cadena de suministro y control durante todo el ciclo de vida. El escaneo en el extremo sigue siendo importante, pero no sustituye la gobernanza de datos aguas arriba, el control de acceso al modelo, los guardrail de la aplicación, las pruebas continuas ni la transparencia. El valor de Safety by Design está en convertir la seguridad de un proyecto puntual en un sistema que se puede mantener.
